% Page 050

\seite{50}

\abschnitt{Spaziergang}
\pstart
\margmark{Am 21. August unte\\gemeinschaftlich e}%
rnahmen die vier Schulen der Pfarrei Seffern\ob
inen Spaziergang nach Hammerles.

\trenner

\pend
\abschnitt{Ernte}
\pstart

\margmark{Die diesjährige Er\\Grumeternte war st\\die trocknen aber\\schlechtes Wetter\\auch Obst, nur in\\so hat wenig auf.}%
nte ist allgemein schlecht gewesen und die\ob
ellenweise ganz gut, manche Wiesen auf\ob
mäßig, weil im Hochsommer lange Zeit\ob
war, wollte muß rings sagen\unsicher. Es gab\ob
den Feldern war die Ernte geringer. Schadfrösten\unsicher\ob
Die Obstpreise waren hoch.

\trenner

\pend
\abschnitt{Schulfestlichkeiten}
\pstart

\margmark{Der Geburtstag Sr.\\27. im Saale zu Se\\im Lokale des Herr\\abwechselnd Lieder\\Ansprache an die K}%
 Majestät des Kaisers wurde wie alljährlich, am\ob
ffern gemeinschaftlich von den Schulen der Pfarrei\ob
n P.\ Schillings dahier gefeiert. Es wurden\ob
 gesungen und Gedichte vorgetragen. Die\ob
inder hielt Pfarrer Worte\unsicher{} aus Kutenbach.

\trenner

\pend
\abschnitt{Kulturfestlichkeiten}
\pstart

\margmark{Am 10. März wurde\\len der Pfarrei ge\\1813 gegen den Fei\\vorgetragen. Die A}%
im Pfarrsaale zu Seffern eine gemeinsame Feier der Schü\ohb
halten, zur Erinnerung an die 100jährige Erhebung Preußens\ob
nd. Es wurden abwechselnd Lieder gesungen und Gedichte\ob
nsprache an die Kinder hielt Pfarrer Kremer aus Seffern.

\margmark{Ges. Seffern 12. 1\\Firmung}%
913        Chronikführung.\ob
wurde am 9.\ Mai durch den Ortsschulinspektor Herrn Pfarrer Kremer aus Sefferweich\unsicher

\pend
\jahresueberschrift{Das Schuljahr 1913-14}
\pstart

\margmark{Entlassen wurden i\\genommen wurden 8\\alle sind katholis}%
n diesem Jahre 5 Kinder Kapitel\unsicher, welche sich der Landwirtschaft auf\ohb
Kinder. Die Schülerzahl beträgt 26 Koven\unsicher{} und 14 Mädchen, 12 Knaben\ob
ch.
\pend
